\section*{अध्याय \ref{ch:b2:ellingham} --- एलिंगम आरेख और धातुकर्म}

\begin{solution}{exo:b2:ellingham:1}
\ce{4Cu + O2 -> 2Cu2O}; \ce{4/3Al + O2 -> 2/3Al2O3}; \ce{2C + O2 -> 2CO};
\ce{C + O2 -> CO2}।
\end{solution}

\begin{solution}{exo:b2:ellingham:2}
$\Delta_r H^\circ = 2(-350.46) = \qty{-700.92}{kJ}$; $\Delta_r S^\circ = 2(43.65)
- 2(41.63) - 205.15 = \qty{-201.1}{J/K}$। अतः $\Delta_r G^\circ = -700.9 +
0.2011\,T$ (\unit{kJ}, $T$ \unit{K} में), जो ज़िंक के गलनांक तक मान्य है।
\end{solution}

\begin{solution}{exo:b2:ellingham:3}
\qty{1500}{K} पर \ce{Cr2O3} की रेखा लगभग \qty{-495}{kJ} पर है। इससे नीचे, अतः
क्रोमियम ऑक्साइड का अपचयन करने में समर्थ: सिलिकॉन, टाइटेनियम, ऐलुमिनियम, मैग्नीशियम,
कैल्सियम। इससे ऊपर, असमर्थ: ताँबा, लोहा, ज़िंक। कार्बन, \qty{-488}{kJ} पर,
ठीक ऊपर है: यह क्रोमियम ऑक्साइड का अपचयन थोड़े ऊँचे ताप पर ही करता है,
और तब कार्बाइड से भरा क्रोमियम देता है --- यही एक कारण है कि
कार्बन-मुक्त मिश्रधातुओं के लिए क्रोमियम ऐलुमिनोतापन से बनाया जाता है।
\end{solution}

\begin{solution}{exo:b2:ellingham:4}
प्रवणता $-\Delta_r S^\circ$ है, जो गैस की मात्रा के परिवर्तन से नियंत्रित होती है।
\ce{C + O2 -> CO2}: गैस का एक मोल एक देता है, $\Delta_r S^\circ \approx 0$,
क्षैतिज। \ce{2C + O2 -> 2CO}: एक से दो, $\Delta_r S^\circ > 0$,
रेखा गिरती है। \ce{2CO + O2 -> 2CO2}: तीन से दो, $\Delta_r S^\circ < 0$,
रेखा चढ़ती है, धातुओं की रेखाओं की तरह।
\end{solution}

\begin{solution}{exo:b2:ellingham:5}
$\Delta_r G^\circ(\qty{600}{K}) = -181.6 + 600 \times 0.2165 = \qty{-51.7}{kJ}$,
$p_{\ce{O2}} = p^\circ\exp(-51\,710/(8.314 \times 600)) = \qty{3.2e-5}{bar}$। वायु में
\qty{0.21}{bar} है, कहीं अधिक: \qty{600}{K} पर पारा ऑक्सीकृत होता है (धीरे-धीरे
--- इसी प्रकार ऑक्साइड पहली बार बनाया गया था), और ऑक्साइड लगभग
\qty{730}{K} से ऊपर ही विघटित होता है।
\end{solution}

\begin{solution}{exo:b2:ellingham:6}
$\Delta_r G^\circ = -1582.3 - (-743.5) = \qty{-838.8}{kJ}$। अभिकर्मक
दो ठोस हैं जो केवल अपने कणों की सतह पर संपर्क में हैं, और अभिक्रिया की
सक्रियण ऊर्जा ऊँची है: उसे स्थानीय रूप से प्रज्वलित करना पड़ता है (मैग्नीशियम का फ़्यूज़);
फिर उसके द्वारा मुक्त ऊष्मा उसे चलाती रहती है।
\end{solution}

\begin{solution}{exo:b2:ellingham:7}
\ce{2C + O2 -> 2CO} की रेखा \ce{2Fe + O2 -> 2FeO} की रेखा को लगभग
\qty{1040}{K} पर काटती है; उसके ऊपर \ce{FeO + C -> Fe + CO} की \omterm{def:b2:reaction-free-energy:reaction-gibbs}{मानक अभिक्रिया
गिब्स ऊर्जा} ऋणात्मक है।
\end{solution}

\begin{solution}{exo:b2:ellingham:8}
$\Delta_r G^\circ = 2(-209.1) - (-396.0) = \qty{-22.2}{kJ}$, $K^\circ =
\exp(22\,200/(8.314 \times 1100)) = 11.3$। तब $x^2/(1 - x) = 11.3$, $x =
0.92$। \qty{700}{K} तक ठंडा करने पर साम्य वापस \ce{CO2} की ओर खिसकता है
($x = 0.016$): कार्बन मोनॉक्साइड असमानुपातित होती है, \ce{2CO -> C + CO2}, और
कालिख जमती है --- धीरे-धीरे, जब तक लोहे जैसी कोई सतह इसे उत्प्रेरित न करे।
\end{solution}

\begin{solution}{exo:b2:ellingham:9}
जल की रेखा (गैस के तीन मोल दो देते हैं) \ce{CO}/\ce{CO2} रेखा से कम
खड़ी उठती है, जो भी तीन से दो है परंतु एन्ट्रॉपी की अधिक हानि के साथ।
दोनों लगभग \qty{1100}{K} पर कटती हैं: नीचे \ce{CO}/\ce{CO2} रेखा नीची है और
कार्बन मोनॉक्साइड प्रबलतर अपचायक है; ऊपर जल की रेखा नीची है और
हाइड्रोजन जीतती है। यह वही कथन है जो उच्च ताप पर जल-गैस साम्य
\ce{CO + H2O <=> CO2 + H2} के बाईं ओर विस्थापन का है।
\end{solution}

\begin{solution}{exo:b2:ellingham:10}
प्रति मोल \ce{O2}, अर्थात् \ce{2MgO + Si -> SiO2 + 2Mg(g)} के लिए: $\Delta_r
G^\circ = -643.7 - (-845.5) = \qty{+201.8}{kJ}$। सक्रियता 1 वाले सिलिकॉन और सिलिका
के साथ $\Delta_r G = \Delta_r G^\circ + RT\ln(p_{\ce{Mg}}/p^\circ)^2$, जो ऋणात्मक है
जब $p_{\ce{Mg}} < p^\circ\exp(-201\,800/(2 \times 8.314 \times 1500)) =
\qty{3e-4}{bar}$। निर्वात भट्ठी मैग्नीशियम वाष्प को बनते ही पंप करके हटा देती है
और उसे ठंडी सतह पर संघनित करती है, इसलिए उसका दाब कभी उस मान तक नहीं पहुँचता;
भार में मिलाया गया चूना सिलिका को बाँधकर उसकी सक्रियता घटाता है,
जो और सहायता करता है।
\end{solution}

\begin{solution}{exo:b2:ellingham:11}
चार स्पीशीज़, एक अभिक्रिया, तीन प्रावस्थाएँ (दो ठोस और गैस): $v = 4 - 1 +
2 - 3 = 2$। दिए गए ताप और दाब पर अनुपात $p_{\ce{CO}}/p_{\ce{CO2}}$
नियत है: गैस अपनी सारी कार्बन मोनॉक्साइड आयरन ऑक्साइड को नहीं दे सकती,
और शीर्ष गैस में अब भी उसकी बड़ी मात्रा रहती है --- जिसे वायु को पहले से गर्म
करने के लिए जलाया जाता है।
\end{solution}

\begin{solution}{exo:b2:ellingham:12}
$\log K = 2(0.34 + 0.41)/0.059 = 25.4$, $K = \num{3e25}$। $[\ce{Fe^2+}] = \qty{1.0}{mol/L}$
के साथ $[\ce{Cu^2+}] = 1/K \approx
\qty{4e-26}{mol/L}$: ताँबा पूरी तरह हट जाता है।
\end{solution}

\begin{solution}{pb:b2:ellingham:1}
\textbf{1.} $\Delta_r H^\circ = \qty{-700.92}{kJ}$, $\Delta_r S^\circ =
\qty{-201.1}{J/K}$।
\textbf{2.} $\Delta_r G^\circ = -700.9 + 0.2011\,T$ (\unit{kJ}),
\qty{692.7}{K} से नीचे।
\textbf{3.} $\Delta_{\text{गलन}}S = 7320/692.7 = \qty{10.6}{J/(K.mol)}$;
$\Delta_{\text{वाष्पन}}S = 115\,300/1180.2 = \qty{97.7}{J/(K.mol)}$।
\textbf{4.} द्रव ज़िंक: $\Delta_r H^\circ = -700.92 - 2(7.32) = \qty{-715.6}{kJ}$,
$\Delta_r S^\circ = -201.1 - 2(10.6) = \qty{-222.2}{J/K}$: $\Delta_r G^\circ =
-715.6 + 0.2222\,T$।
\textbf{5.} गैसीय ज़िंक: $\Delta_r H^\circ = -715.6 - 2(115.3) = \qty{-946.2}{kJ}$,
$\Delta_r S^\circ = -222.2 - 2(97.7) = \qty{-417.7}{J/K}$: $\Delta_r G^\circ =
-946.2 + 0.4177\,T$।
\textbf{6.} प्रवणताएँ $0.201$, $0.222$ और $0.418$ \unit{kJ/K}। क्वथनांक से ऊपर
अभिक्रिया गैस का एक मोल नहीं, बल्कि तीन मोल उपभुक्त करती है: एन्ट्रॉपी की
हानि, और इसलिए प्रवणता, दुगुनी हो जाती है।
\textbf{7.} \qty{692.7}{K} पर: $-700.9 + 139.3 = -715.6 + 153.9 = \qty{-561.6}{kJ}$;
\qty{1180.2}{K} पर: $-715.6 + 262.3 = -946.2 + 492.9 = \qty{-453.3}{kJ}$।
रेखाएँ जुड़ती हैं, जैसा उन्हें जुड़ना चाहिए।
\textbf{8.} प्रवणता $(-453.0 + 418.2)/200 = \qty{-0.174}{kJ/K}$, $a = -418.2 +
0.174 \times 1100 = \qty{-226.8}{kJ}$: $\Delta_r G^\circ = -226.8 - 0.174\,T$।
\textbf{9.} $\Delta_r S^\circ = \qty{+174}{J/K}$: गैस का एक मोल दो देता है।
\textbf{10.} \ce{2ZnO + 2C -> 2Zn(g) + 2CO}, $\Delta_r G^\circ = (-226.8 - 0.174\,T)
- (-946.2 + 0.4177\,T) = 719.4 - 0.592\,T$ (\unit{kJ})।
\textbf{11.} $T = 719.4/0.592 = \qty{1215}{K}$ पर शून्य।
\textbf{12.} अपने क्वथनांक से ऊपर ज़िंक वाष्प के रूप में बनता है: अभिक्रिया
ठोसों से गैस बनाती है, उसकी एन्ट्रॉपी बड़ी और धनात्मक है, और वाष्प
भार को छोड़ देती है, जो आसवन द्वारा धातु को अयस्क और कोक से अलग करता है।
\textbf{13.} चार स्पीशीज़, एक अभिक्रिया, एक विशेष शर्त ($p_{\ce{Zn}} =
p_{\ce{CO}}$), तीन प्रावस्थाएँ: $v = 4 - 1 - 1 + 2 - 3 = 1$। दाब नियत करने से
साम्य का ताप नियत हो जाता है।
\textbf{14.} प्रति मोल कार्बन मोनॉक्साइड ज़िंक वाष्प का एक मोल:
$p_{\ce{Zn}} = p_{\ce{CO}} = \qty{0.5}{bar}$।
\textbf{15.} प्रति मोल \ce{ZnO}, $\Delta_r G^\circ = \frac12(719.4 - 0.592 \times
1300) = \qty{-25.1}{kJ/mol}$, $K^\circ = \exp(25\,100/(8.314 \times 1300)) = 10$;
$Q = 0.25 < K^\circ$: अपचयन चलता है।
\textbf{16.} प्रतिच्छेद से ठीक ऊपर अपचयन पहले से ही अनुकूल है (हर गैस के
\qty{0.5}{bar} के साथ, लगभग \qty{1170}{K} से); अधिक गर्म करना ईंधन माँगता है,
मिट्टी के भभकों को घिसता है और अधिक ज़िंक नहीं बनाता, क्योंकि अभिक्रिया वैसे भी
पूर्णता तक जाती है।
\textbf{17.} \ce{Zn(g) + CO2 -> ZnO + CO}।
\textbf{18.} ज़िंक ऑक्साइड की रेखा \ce{CO}/\ce{CO2} रेखा से नीचे है
(\qty{1200}{K} पर $-445$, जबकि \qty{-357}{kJ}): ज़िंक वाष्प कार्बन
डाइऑक्साइड का अपचयन करके फिर ऑक्साइड बन जाती है। भभके में गर्म कार्बन किसी भी
\ce{CO2} को नष्ट कर देता है (बूदुआर); ठंडे संघनित्र में कोई नहीं करता, और \ce{CO2} --- भीतर
रिसती वायु से, या ठंडा होने पर \ce{2CO -> C + CO2} से --- ज़िंक को एक
धूसर ऑक्सीकृत चूर्ण में बिगाड़ देती है।
\textbf{19.} धीमा शीतलन वाष्प को लंबे समय तक उस परिसर में छोड़ता है जहाँ
वह फिर ऑक्सीकृत होती है, और द्रव धातु के स्थान पर बारीक धूल देता है; द्रव तक
शीघ्र संघनन दोनों को सीमित करता है।
\textbf{20.} ज़िंक के $10^6/65.38 = \qty{1.53e4}{mol}$, कार्बन के उतने ही:
$1.53 \times 10^4 \times 12.011 = \qty{184}{kg}$ (व्यवहार में कहीं अधिक, भभकों को
गर्म करने के लिए)।
\textbf{21.} प्रति मोल ज़िंक $\Delta_r H^\circ = \frac12(-226.8 + 946.2) = \qty{+360}{kJ}$;
प्रति टन, $1.53 \times 10^4 \times 360 = \qty{5.5e6}{kJ}$, अर्थात्
\qty{5.5}{GJ} --- जो भभकों के चारों ओर और कोयला जलाकर दी जाती है।
\textbf{22.} मानक परिस्थितियों में कार्बन ज़िंक ऑक्साइड का अपचयन
$\boldsymbol{T \approx \qty{1.22e3}{K}}$ से ऊपर करता है।
\end{solution}
